% The weak reaction that swaps neutrons and protons before BBN, n nu_e <-> p e-,
% drawn twice: with the nucleons, then with their quarks (udd -> udu: one d -> u).
% (The other one, n e+ <-> p nubar_e, is the same diagram with the lepton line crossed.)
% Time runs upwards, as in nu-e-reactions.tex and fermi-blob.tex. Uses feynmp via feynmp-auto.
% Build (in this directory):
%   pdflatex -shell-escape n-p-weak.tex && pdflatex -shell-escape n-p-weak.tex
%   pdftocairo -svg n-p-weak.pdf ../n-p-weak.svg
\documentclass[border=6pt]{standalone}
\usepackage{feynmp-auto}
\usepackage{xcolor}

\newif\ifdark \darktrue          % \darkfalse = black on white
\newcommand{\figfont}{\LARGE}    % font size for labels and captions
\newcommand{\diagramgap}{\hspace{24mm}}  % space between the two diagrams

\ifdark \pagecolor{black} \fi
\unitlength=1mm

\begin{document}
\ifdark \color{white} \fi
\figfont
\begin{fmffile}{npweak}
  \ifdark \fmfcmd{foreground := white; background := black;} \fi
  \fmfset{thin}{0.8pt}
  \fmfset{arrow_len}{3mm}
  \fmfset{wiggly_len}{3mm}
  \begin{tabular}{c@{\diagramgap}c}
    % --- nucleons: n nu_e -> p e- by W exchange ---
    $\vcenter{\hbox{\fmfframe(4,8)(4,8){% margins (left,top)(right,bottom) leave room for labels
    \begin{fmfgraph*}(40,40)
      \fmfbottom{n,nue}
      \fmftop{p,e}
      \fmf{fermion}{n,v1,p}
      \fmf{fermion}{nue,v2,e}
      \fmf{boson,label=$W$,label.side=left}{v1,v2}
      \fmflabel{$n$}{n}
      \fmflabel{$p$}{p}
      \fmflabel{$\nu_e$}{nue}
      \fmflabel{$e^-$}{e}
    \end{fmfgraph*}}}}$
    &
    % --- quarks: two spectators, one d -> u ---
    $\vcenter{\hbox{\fmfframe(4,8)(4,8){%
    \begin{fmfgraph*}(50,40)
      \fmfstraight
      \fmfbottom{u1i,d1i,d2i,nue}
      \fmftop{u1o,d1o,u2o,e}
      % three parallel quark lines, each bent like the n -> p line on the left;
      % the outer one (d -> u) carries the W vertex, the inner two bend at plain vertices k1, k2
      \fmfforce{(0,0)}{u1i}       \fmfforce{(.2w,.5h)}{k1}    \fmfforce{(0,h)}{u1o}
      \fmfforce{(.13w,0)}{d1i}    \fmfforce{(.33w,.5h)}{k2}   \fmfforce{(.13w,h)}{d1o}
      \fmfforce{(.26w,0)}{d2i}    \fmfforce{(.46w,.5h)}{v1}   \fmfforce{(.26w,h)}{u2o}
      \fmfforce{(w,0)}{nue}       \fmfforce{(w,h)}{e}
      \fmfforce{(.8w,.5h)}{v2}
      \fmf{fermion}{u1i,k1,u1o}
      \fmf{fermion}{d1i,k2,d1o}
      \fmf{fermion}{d2i,v1,u2o}
      \fmf{fermion}{nue,v2,e}
      \fmf{boson,label=$W$,label.side=left}{v1,v2}
      \fmflabel{$u\vphantom{d}$}{u1i} \fmflabel{$d$}{d1i} \fmflabel{$d$}{d2i}  % vphantom: same baseline for u and d
      \fmflabel{$u\vphantom{d}$}{u1o} \fmflabel{$d$}{d1o} \fmflabel{$u\vphantom{d}$}{u2o}
      \fmflabel{$\nu_e$}{nue}
      \fmflabel{$e^-$}{e}
    \end{fmfgraph*}}}}$
    \\[2mm]
    $n\,\nu_e \leftrightarrow p\,e^-$ & $udd \to udu$
  \end{tabular}
\end{fmffile}
\end{document}
